\section{Recovery}
\label{sec:recovery}

Recovery restores a node to a correct state after a restart or crash. \vannak{}
relies on three durable substrates: append-only segments, file-backed \raft{}
storage, and \raft{}-replicated checkpoint manifests.

\subsection{Checkpoint Manifest}

A \texttt{CheckpointManifest} is a shard-level recovery checkpoint replicated
through \raft{}. It records, per segment, the last consumed record offset, so a
recovering node can skip already-processed data.

\begin{lstlisting}
pub struct CheckpointManifest {
    pub checkpoint_id: String,
    pub node_id: NodeId,
    pub epoch: CheckpointEpoch,
    pub segment_offsets: BTreeMap<SegmentId, RecordOffset>,
    pub metadata_version: Option<String>,
    pub checksum: u64,
}
\end{lstlisting}

The helper \texttt{latest\_offset\_for} returns the recorded offset for a given
segment id, which seeds checkpoint-aware replay.

\subsection{Recovery Flow}

On restart a node reconstructs its operational state in order:

\begin{enumerate}
\item Load the \raft{} log and persistent state from disk.
\item Install the latest \texttt{ClusterStateMachine} snapshot and replay any
  log entries after it.
\item Read the relevant \texttt{CheckpointManifest} and
  \texttt{MetadataOutboxCheckpoint} records from the control state.
\item Replay each owned segment after its checkpoint offset to rebuild pending
  outbox state.
\item Resume delivery of pending entries through the \texttt{IptoWriter}.
\end{enumerate}

\subsection{Outbox Recovery}

Outbox recovery uses \texttt{replay\_metadata\_outbox\_segment\_after} with the
checkpoint offset from \texttt{MetadataOutboxCheckpoint}. Records at or before
the offset were already acknowledged and are skipped; records after it are
re-queued as pending. The returned \texttt{MetadataOutboxReplaySummary} confirms
how many records were scanned, skipped, and replayed.

\texttt{SegmentBackedMetadataOutbox::recover\_after} combines that replay with
\texttt{SegmentWriter::open\_append}. Recovery validates the existing segment,
reconstructs pending outbox entries after the checkpoint, preserves the segment
manifest counters, and reopens the same segment for subsequent durable appends.
\texttt{VannakService::recover\_after} wraps this as the single-node service
recovery constructor.

\begin{vnote}
Because the durable outbox syncs each payload to its segment before exposing it
as pending, no acknowledged write can be lost across a crash, and no pending
write is forgotten: replay rebuilds it from the segment. A recovered outbox can
then continue appending new metadata writes to the validated segment.
\end{vnote}

\subsection{Shard-Range Replay}

When recovery coincides with a placement change, the shard-range replay
primitive (Section~\ref{sec:rebalancing}) selects only the records for a shard
range, allowing a recovering node to re-home a subset of pending entries to a
new \ipto{} target rather than the original one.

\subsection{Corruption Detection}

Segments are self-describing and checksummed.
\texttt{SegmentReader::open} rejects a file whose header is not the
\texttt{VANNAK01} magic with \texttt{SegmentError::InvalidMagic}.
\texttt{read\_next} recomputes each record's checksum and returns
\texttt{SegmentError::ChecksumMismatch} (with the offset and both checksums) when
they differ, so corruption is detected at the exact record during replay rather
than silently propagated.

\subsection{State Machine Snapshot}

\texttt{ClusterStateMachine::snapshot} serializes the entire control state to
JSON, and \texttt{restore} rebuilds it by replaying the snapshot through the
normal command path. Because restore reuses \texttt{apply}, an installed
snapshot is subject to the same validation as live commands, which guards
against installing an inconsistent snapshot.

\subsection{Raft Wrapping}

The \texttt{vannak-node} binary wraps the control state in \raft{} using a
\texttt{FileLogStore} and a \texttt{FilePersistentStateStore}. These persist the
\raft{} log and the persistent node state under the data directory, so the
\raft{} layer itself survives a restart and resumes from its committed index.

\subsection{File-Backed Persistence}

Each node is given a data directory and persists two subdirectories:

\begin{table}[htbp]
\centering
\begin{tabular}{ll}
\hline
Path & Contents \\
\hline
\filepath{<data-dir>/log} & \raft{} log entries (FileLogStore) \\
\filepath{<data-dir>/state} & persistent \raft{} state (FilePersistentStateStore) \\
\hline
\end{tabular}
\caption{Per-node file-backed persistence layout.}
\end{table}

In the cluster Compose file each node mounts a named volume at \filepath{/data},
so this state survives container restarts.

\subsection{Crash Recovery}

After an unclean stop, a node restarts by loading its file-backed \raft{} state,
re-electing or rejoining a leader, installing the latest state machine snapshot,
and replaying committed log entries. Operational pending state is then rebuilt
from segments using the checkpoints, as described above. No special operator
action is required for a single-node crash within a healthy cluster.

\subsection{Graceful Shutdown}

A node runs an accept loop over its transport listener; shutdown is process
termination, typically via \shellcmd{docker compose -f docker-compose.cluster.yml down}.
This is safe because all durable state is already on disk: outbox segments are
synced on enqueue, sealed segment manifests are recorded through \raft{}, and
the \raft{} log and state are file-backed. A restarted node recovers from these.

\subsection{Degraded Mode}

\vannak{} degrades rather than failing hard in several situations:

\begin{itemize}
\item If DNS SRV discovery resolves no peers, a node logs a warning and starts
  as a single node.
\item During reconfiguration, reads use \texttt{resolve\_with\_fallback} and,
  as a last resort, broadcast across \texttt{all\_ipto\_instances}.
\item A permanent \texttt{IptoWriteError} leaves the affected outbox entry in
  the \texttt{Failed} state for operator attention, while other entries continue
  to deliver.
\end{itemize}

\endinput
